\documentclass[12pt]{article}
\usepackage{amsmath, amssymb}

\begin{document}

\title{CSE400 -- Lecture 18 Scribe}
\date{}
\maketitle

\section*{1. Definitions and Notations}

\subsection*{a) Definition}

\textbf{Marginal PMF:} \\
The marginal PMF of one variable is obtained by summing the joint PMF over all possible values of the other variable.

\textbf{Correlation (Product Moment):}
\[
R_{XY} = E[XY]
\]
It represents the average value of the product $XY$.

\textbf{Covariance:}
\[
\mathrm{Cov}(X,Y) = E[(X-\mu_X)(Y-\mu_Y)]
\]
It captures whether deviations from means occur in the same or opposite direction.

\textbf{Pearson’s Correlation Coefficient:}
\[
\rho_{XY} = \frac{\mathrm{Cov}(X,Y)}{\sigma_X \sigma_Y}
\]
It measures the strength and direction of linear relationship.

\textbf{Jointly Gaussian Random Variables:} \\
A pair $(X,Y)$ is jointly Gaussian if its joint density has a bivariate Gaussian form and the marginals are Gaussian.

\subsection*{b) Notation}

\[
p_{XY}(x,y) = P(X=x,Y=y)
\]

\[
p_X(x) = \sum_y p_{XY}(x,y), \quad p_Y(y) = \sum_x p_{XY}(x,y)
\]

\[
R_{XY} = E[XY]
\]

\[
\mu_X = E[X], \quad \mu_Y = E[Y]
\]

\[
\mathrm{Cov}(X,Y) = E[XY] - E[X]E[Y]
\]

\[
\rho_{XY} = \frac{\mathrm{Cov}(X,Y)}{\sqrt{\mathrm{Var}(X)\mathrm{Var}(Y)}}
\]

\textbf{Joint Gaussian PDF:}

\[
f_{XY}(x,y) =
\frac{1}{2\pi \sigma_X \sigma_Y \sqrt{1-\rho^2}}
\exp\left(
-\frac{1}{2(1-\rho^2)}
\left[
\left(\frac{x-\mu_X}{\sigma_X}\right)^2 +
\left(\frac{y-\mu_Y}{\sigma_Y}\right)^2 -
2\rho \left(\frac{x-\mu_X}{\sigma_X}\right)\left(\frac{y-\mu_Y}{\sigma_Y}\right)
\right]
\right)
\]

\section*{2. Assumptions and Preliminaries}

\begin{itemize}
\item $X$ and $Y$ are discrete or continuous random variables with valid joint distributions.
\item Total probability sums to 1.
\item Expectations $E[X], E[Y], E[XY]$ exist.
\item Covariance uses centered variables (subtracting means).
\item Correlation depends on scale and is computed about origin.
\item Zero covariance implies uncorrelated but does not imply independence.
\end{itemize}

\section*{3. Main Result}

\subsection*{a) Proof}

\textbf{Marginal PMF derivation:}
\[
p_X(x) = \sum_y p_{XY}(x,y), \quad p_Y(y) = \sum_x p_{XY}(x,y)
\]

\textbf{Covariance derivation:}
\[
\mathrm{Cov}(X,Y) = E[(X-\mu_X)(Y-\mu_Y)]
\]
\[
= E[XY - X\mu_Y - Y\mu_X + \mu_X \mu_Y]
\]
\[
= E[XY] - \mu_Y E[X] - \mu_X E[Y] + \mu_X \mu_Y
\]
\[
= E[XY] - E[X]E[Y]
\]

\textbf{Properties:}

\begin{itemize}
\item Symmetry:
\[
\mathrm{Cov}(X,Y) = \mathrm{Cov}(Y,X)
\]

\item Variance:
\[
\mathrm{Cov}(X,X) = \mathrm{Var}(X)
\]

\item Linearity:
\[
\mathrm{Cov}(aX+b, cY+d) = ac \, \mathrm{Cov}(X,Y)
\]

\item Independence:
\[
X \perp Y \Rightarrow \mathrm{Cov}(X,Y)=0
\]

\item Pearson Coefficient:
\[
\rho_{XY} = \frac{\mathrm{Cov}(X,Y)}{\sigma_X \sigma_Y}, \quad -1 \leq \rho_{XY} \leq 1
\]
\end{itemize}

\subsection*{b) Conclusion}

\begin{itemize}
\item Marginal PMFs are obtained by summing over the other variable.
\item Correlation $E[XY]$ mixes mean, spread, and dependence.
\item Covariance measures centered dependence.
\item Pearson coefficient normalizes covariance and gives strength and direction.
\item Joint Gaussian variables follow a bivariate normal distribution with parameters $\mu_X, \mu_Y, \sigma_X, \sigma_Y, \rho$.
\end{itemize}

\section*{4. Working Examples}

\textbf{Example 1: Marginal PMF}

\[
\begin{array}{c|cc}
 & y=0 & y=1 \\
\hline
x=0 & \frac{1}{8} & \frac{1}{4} \\
x=1 & \frac{3}{8} & \frac{1}{4}
\end{array}
\]

\textbf{Steps:}

\[
p_X(0) = \frac{1}{8} + \frac{1}{4} = \frac{3}{8}, \quad
p_X(1) = \frac{3}{8} + \frac{1}{4} = \frac{5}{8}
\]

\[
p_Y(0) = \frac{1}{8} + \frac{3}{8} = \frac{1}{2}, \quad
p_Y(1) = \frac{1}{4} + \frac{1}{4} = \frac{1}{2}
\]

\textbf{Final:}
\[
p_X = \{3/8, 5/8\}, \quad p_Y = \{1/2, 1/2\}
\]

\vspace{10pt}

\textbf{Example 2: Covariance}

Given:
\[
Y = -2X + 3
\]

\[
XY = -2X^2 + 3X
\]

\[
E[Y] = -2E[X] + 3
\]

\[
\mathrm{Cov}(X,Y) = -2E[X^2] + 2(E[X])^2
\]

\[
\mathrm{Var}(X) = E[X^2] - (E[X])^2
\]

\textbf{Final:}
\[
\mathrm{Cov}(X,Y) = -2 \, \mathrm{Var}(X)
\]

\vspace{10pt}

\textbf{Example 3: Continuous Case}

\[
f_{XY}(x,y) = \frac{xy}{96}, \quad 0<x<4,\; 1<y<5
\]

\[
E[X] = \frac{8}{3}, \quad E[Y] = \frac{31}{9}
\]

\[
E[XY] = \frac{248}{27}
\]

\[
E[2X+3Y] = \frac{47}{3}
\]

\[
\mathrm{Var}(X) = \frac{8}{9}, \quad \mathrm{Var}(Y) = \frac{92}{81}
\]

\[
\mathrm{Cov}(X,Y) = 0, \quad \rho_{XY} = 0
\]

\vspace{10pt}

\textbf{Example 4: Joint Gaussian Interpretation}

$X$: CO$_2$ concentration \\
$Y$: Temperature anomaly

Joint behavior approximated by Gaussian distribution due to multiple random factors.

If $\rho > 0$, higher CO$_2$ corresponds to higher temperature.

\end{document}